% ============================================================================
% Fichier  : partie4/P04_C11_gerer_vulnerabilites.tex
% Titre    : Gérer les Vulnérabilités
% Auteur   : MINKA MI NGUIDJOÏ Thierry Emmanuel
% Version  : 1.0 -- novembre 2026
% Corrigés : corriges/Corr_C11.tex
% ============================================================================
% Sources : NIST SP 800-40 r4 (gestion des correctifs comme maintenance
% préventive) ; ISO/IEC 30111 et 29147 (traitement et divulgation côté
% éditeur) ; ISO/IEC 27002:2022, mesure 8.8 (gestion des vulnérabilités
% techniques, paraphrasée) ; catalogue KEV ; EPSS.
% Faits nouveaux (bible) : politique SOLITRANS de juin 2025, classes P1 à P4
% et délais 7/15/30/90 jours ; suivi T3 2025 à T1 2026 ; exceptions E-01 à
% E-04 ; NYANGA adopte la méthode en septembre 2025.
% ============================================================================

\chapter{Gérer les Vulnérabilités}
\label{chap:gerer}
\chaptermark{Gérer les vulnérabilités}

\epigraph{%
    « Une vulnérabilité corrigée est une vulnérabilité\\
    dont on a vérifié la correction. »%
}{Principe du cours}

\niveauChapitre{Intermédiaire}{%
    Maîtrise des Chapitres~\ref{chap:vulnerabilite} et~\ref{chap:identifier}
}

\begin{boxobjectifs}
\begin{itemize}
    \item Décrire le cycle de gestion des vulnérabilités, de
          l'identification à la vérification.
    \item Rédiger une politique de priorité qui combine exploitation,
          exposition et importance de l'actif, avec des délais.
    \item Choisir entre correctif, changement de configuration et mesure
          compensatoire.
    \item Gérer une exception~: propriétaire, durée, mesure compensatoire,
          revue.
    \item Construire et lire les indicateurs d'un programme.
\end{itemize}
\end{boxobjectifs}

\bigskip

% ============================================================================
\section*{Avant de commencer : quarante failles par semaine}
% ============================================================================

Douala, juin 2025. Mme~EBODE présente sa politique au comité de direction
de SOLITRANS. Elle ne propose pas de corriger plus~: l'équipe système de
M.~FOUDA traite environ quarante failles par semaine, et rien ne permettra
de doubler ce rythme avant l'an prochain. Elle propose de corriger dans le
bon ordre, avec des délais écrits, et de vérifier ce qui est fait. Le
directeur financier demande combien de failles resteront ouvertes à la fin
de l'année. Elle répond qu'il en restera beaucoup, et qu'elle saura dire
lesquelles, pourquoi, et qui en a accepté le risque.

\begin{boxalert}
\textbf{Violation typique.} Déclarer une faille corrigée sur la foi du
ticket fermé par l'équipe système, sans nouvelle analyse de l'actif.

\smallskip
\textbf{Règle~:} une correction se vérifie par un nouveau balayage, ou par
une lecture directe de la version installée. Le ticket fermé est une
déclaration~; le balayage qui ne trouve plus la faille est une constatation.
\end{boxalert}

\separateur

% ============================================================================
\section{Le Cycle de Gestion}
\label{sec:c11-cycle}
% ============================================================================

Le guide NIST SP~800-40~\cite{nist80040} présente la gestion des correctifs
comme une maintenance préventive, planifiée, et non comme une suite
d'urgences. La mesure 8.8 de l'ISO/IEC~27002:2022~\cite{iso27002_2022},
consacrée à la gestion des vulnérabilités techniques, demande d'obtenir
l'information sur les vulnérabilités, d'évaluer l'exposition de
l'organisation et de prendre les mesures appropriées (paraphrase).

\begin{center}
\begin{tikzpicture}[
    etape/.style = {draw=CouleurBordDef, fill=CouleurDef, rounded corners=3pt,
        minimum width=2.4cm, minimum height=0.9cm, align=center,
        font=\small\bfseries\color{CouleurPrimaire}},
    arrow/.style = {-{Stealth[length=4pt]}, thick, color=CouleurSecondaire}
]
\node[etape] (a) at (0,0)     {Inventorier};
\node[etape] (b) at (3.0,0)   {Identifier};
\node[etape] (c) at (6.0,0)   {Prioriser};
\node[etape] (d) at (9.0,0)   {Traiter};
\node[etape] (e) at (9.0,-1.6) {Vérifier};
\node[etape] (f) at (4.5,-1.6) {Rendre compte};
\draw[arrow] (a) -- (b); \draw[arrow] (b) -- (c); \draw[arrow] (c) -- (d);
\draw[arrow] (d) -- (e); \draw[arrow] (e) -- (f); \draw[arrow] (f) -| (a);
\end{tikzpicture}
\end{center}
\begin{center}{\small\itshape Figure~11.1. Le cycle de gestion des vulnérabilités}\end{center}

Traiter ne veut pas toujours dire corriger. Quatre voies existent.

\begin{tableaucours}{L{2.8cm}L{5.0cm}L{4.4cm}}{%
    \tableauHeader{Voie} &
    \tableauHeader{Exemple} &
    \tableauHeader{Condition}
}
Correctif & Mise à jour de l'éditeur & Testé, appliqué dans une fenêtre,
retour arrière prévu \\
Configuration & Désactiver le service vulnérable, changer un réglage &
Fonction non utilisée, ou réglage documenté par l'éditeur \\
Mesure compensatoire & Filtrage réseau, restriction d'accès, règle de
pare-feu applicatif & Temporaire, en attendant le correctif \\
Acceptation & Exception formelle & Propriétaire, durée, revue (section~11.4) \\
\end{tableaucours}
\vspace{-0.8em}
\begin{center}{\small\itshape Tableau~11.1. Les quatre voies de traitement}\end{center}

% ============================================================================
\section{La Politique de Priorité}
\label{sec:c11-politique}
% ============================================================================

Une politique de priorité traduit en règles écrites ce que le
Chapitre~\ref{chap:vulnerabilite} a montré~: l'exploitation et l'exposition
comptent davantage que le score seul.

\begin{boxscenario}{Opération SOLITRANS. La politique de juin 2025}
\begin{tableaucours}{cL{6.6cm}cL{2.6cm}}{%
    \tableauHeader{Classe} &
    \tableauHeader{Règle d'entrée} &
    \tableauHeader{Délai} &
    \tableauHeader{Décideur d'exception}
}
P1 & Inscrite au catalogue KEV, ou EPSS $\geq 0{,}5$, sur un actif exposé à
Internet ou critique & 7 jours & Directeur général \\
P2 & Inscrite au catalogue KEV sur un actif interne, ou CVSS-B $\geq 9{,}0$
sur un actif exposé & 15 jours & RSSI \\
P3 & CVSS-B $\geq 7{,}0$, hors P1 et P2 & 30 jours & RSSI \\
P4 & Autres constats confirmés & 90 jours & Responsable des
vulnérabilités \\
\end{tableaucours}

\medskip
\textbf{Constatations.} Appliquée aux 2\,947 constats de mai, la politique
place 38 constats en P1, 211 en P2, 904 en P3 et 1\,794 en P4. À quarante
corrections par semaine, les 249 constats P1 et P2 demandent un peu plus de
six semaines, bien au-delà du délai de 15~jours fixé pour les P2.

\textbf{Analyse} (confiance élevée). La politique est tenable pour les P1,
pas encore pour le stock initial de P2~: Mme~EBODE demande un renfort
temporaire de deux semaines, et la fixation du délai à partir de la date
d'entrée en vigueur pour le stock existant.
\end{boxscenario}

\begin{boxinfo}
\textbf{Les délais courent depuis quand~?} Le délai d'une faille se compte à
partir du moment où l'organisation pouvait la connaître~: la publication du
correctif, ou la première détection sur son parc, selon la règle choisie.
Une politique qui ne le précise pas permet de repousser indéfiniment
l'échéance.
\end{boxinfo}

% ============================================================================
\section{Corriger en Production}
\label{sec:c11-production}
% ============================================================================

Un correctif peut casser une application. C'est la raison la plus souvent
donnée pour ne pas l'appliquer, et elle est parfois exacte. La réponse n'est
pas de renoncer, mais d'organiser~: un environnement de test qui reproduit
la production pour les applications critiques, des fenêtres de maintenance
connues des métiers, une procédure de retour arrière éprouvée, et un ordre
de déploiement par vagues (d'abord quelques machines représentatives, puis
le reste).

\begin{boxinfo}
\textbf{Déployer un correctif de classe P1}
\begin{enumerate}
    \item Lire le bulletin~: versions touchées, prérequis, effets connus.
    \item Appliquer une mesure compensatoire immédiate si le correctif
          demande plus de 24~heures.
    \item Tester sur une machine représentative, avec le propriétaire
          métier.
    \item Déployer par vagues, en surveillant les services.
    \item Vérifier par un nouveau balayage authentifié de chaque actif.
    \item Lever la mesure compensatoire, et clore le constat avec la
          preuve.
\end{enumerate}
\end{boxinfo}

% ============================================================================
\section{Les Exceptions}
\label{sec:c11-exceptions}
% ============================================================================

Certaines failles ne peuvent pas être corrigées dans le délai~: le
correctif n'existe pas, l'éditeur ne maintient plus le produit, ou
l'application métier ne fonctionne pas avec la version corrigée. L'exception
transforme ce fait en décision.

\begin{boxdef}
\textbf{\termecle{Exception}}~: décision écrite par laquelle un responsable
habilité accepte, pour une durée limitée, qu'une vulnérabilité reste
ouverte au-delà du délai, avec une mesure compensatoire et une date de
revue.
\end{boxdef}

\begin{boxscenario}{Opération SOLITRANS. Quatre exceptions}
\begin{tableaucours}{cL{3.6cm}L{3.4cm}cL{1.8cm}}{%
    \tableauHeader{Réf.} &
    \tableauHeader{Faille et actif} &
    \tableauHeader{Mesure compensatoire} &
    \tableauHeader{Durée} &
    \tableauHeader{Décideur}
}
E-01 & P2, logiciel de dédouanement en fin de support & Accès limité à
trois postes, réseau isolé & 6 mois & RSSI \\
E-02 & P3, terminaux de pesage, aucun correctif & Filtrage strict, aucun
accès Internet & 12 mois & RSSI \\
E-03 & P1, portail des transitaires, correctif incompatible & Règle de
pare-feu applicatif & 30 jours & Directeur général \\
E-04 & P3, serveur d'archives & Aucune & Sans limite & Non signée \\
\end{tableaucours}

\medskip
\textbf{Analyse} (confiance élevée). E-01 à E-03 sont des exceptions. E-04
n'en est pas une~: sans mesure, sans durée et sans signature, c'est une
faille ouverte que personne n'assume. Elle redevient un constat en retard.
\end{boxscenario}

% ============================================================================
\section{Mesurer le Programme}
\label{sec:c11-mesurer}
% ============================================================================

\begin{tableaucours}{L{4.0cm}L{5.2cm}L{3.0cm}}{%
    \tableauHeader{Indicateur} &
    \tableauHeader{Définition} &
    \tableauHeader{Piège}
}
Délai de correction & Durée médiane entre détection et correction
vérifiée, par classe & La moyenne masque les cas extrêmes \\
Respect des délais & Part des constats corrigés dans le délai de leur classe
& À calculer sur les constats échus \\
Couverture & Part des actifs existants analysés avec authentification &
Souvent omise \\
KEV ouvertes & Nombre de failles du catalogue encore ouvertes & Le seul
indicateur à suivre chaque semaine \\
Âge du stock & Âge des constats ouverts, par tranche & Le stock vieillit
sans bruit \\
\end{tableaucours}
\vspace{-0.8em}
\begin{center}{\small\itshape Tableau~11.2. Les indicateurs d'un programme}\end{center}

\begin{boxlecture}{SOLITRANS, trois trimestres}
\begin{tableaucours}{L{5.6cm}ccc}{%
    \tableauHeader{Indicateur} &
    \tableauHeader{T3 2025} &
    \tableauHeader{T4 2025} &
    \tableauHeader{T1 2026}
}
Couverture authentifiée & 71~\% & 84~\% & 93~\% \\
Délai médian P1 (jours) & 11 & 6 & 5 \\
P1 corrigées dans le délai & 58~\% & 81~\% & 92~\% \\
P2 corrigées dans le délai & 34~\% & 63~\% & 77~\% \\
KEV ouvertes en fin de trimestre & 14 & 5 & 2 \\
Exceptions en cours, toutes signées & \NON{} & \OUI{} & \OUI{} \\
\end{tableaucours}

\medskip
\textbf{Lecture.} Le programme progresse sur toutes les lignes. La
couverture passe de 71 à 93~\%~: une partie des nouvelles failles du T4
vient simplement de ce que l'on regarde mieux, et la hausse apparente du
stock ne traduit pas une dégradation. Les deux KEV ouvertes au T1 2026 sont
les terminaux de pesage, sous l'exception E-02.
\end{boxlecture}

\begin{boxscenario}{Opération NYANGA. Septembre 2025}
\textbf{Constatations.} Six mois après l'incident, NYANGA CRÉDIT adopte la
méthode de SOLITRANS~: inventaire croisé (la passerelle et quatorze autres
équipements y entrent), balayage authentifié mensuel, politique en quatre
classes, suivi hebdomadaire des KEV. Mme~NDOUMBE devient aussi responsable
du programme.

\textbf{Analyse} (confiance moyenne). Avec cette politique, VULN-2024-4471
aurait été classée P1 le 19~décembre 2024, avec une échéance au
26~décembre, soit plus de deux mois avant le premier accès.
\end{boxscenario}

% ============================================================================
\section{Ce que la Gestion demande à chaque Métier}
\label{sec:c11-metiers}
% ============================================================================

\begin{boxmetier}{Analyste en gestion des vulnérabilités}
\textbf{Tenir le programme}
\begin{enumerate}
    \item Classer chaque constat confirmé selon la politique écrite.
    \item Suivre chaque semaine les failles du catalogue KEV ouvertes.
    \item Vérifier chaque correction par un nouveau balayage.
    \item Instruire les exceptions et rappeler leurs dates de revue.
    \item Publier les indicateurs avec la couverture.
\end{enumerate}
\end{boxmetier}

\begin{boxmetier}{RSSI}
Faire approuver la politique et les délais par la direction, et signer les
exceptions de son niveau. Refuser une exception sans mesure compensatoire ni
durée~: elle ne protège personne, sauf de la responsabilité.
\end{boxmetier}

\begin{boxmetier}{Analyste en réponse à incident}
Après un incident, vérifier que la faille exploitée est corrigée partout où
elle existe, et pas seulement sur la machine compromise.
\end{boxmetier}

\separateur

\begin{boxresume}
\textbf{Ce qu'il faut retenir de ce chapitre~:}
\begin{itemize}
    \item Le cycle va de l'inventaire à la vérification~; traiter peut
          vouloir dire corriger, reconfigurer, compenser ou accepter.
    \item Une politique écrite classe les failles par exploitation,
          exposition et importance, avec des délais datés.
    \item Une exception a un décideur, une mesure compensatoire, une durée
          et une revue~; sinon, c'est un retard.
    \item Une correction se vérifie par un nouveau balayage~; un ticket
          fermé ne suffit pas.
    \item Les indicateurs se lisent avec la couverture~: une hausse des
          failles peut traduire un meilleur regard.
\end{itemize}
\end{boxresume}

\begin{boxinfo}
\textbf{Points de contrôle.} Vous avez compris ce chapitre si vous savez~:
\begin{itemize}
    \item classer une faille dans une politique en quatre classes~;
    \item distinguer une exception d'un retard~;
    \item interpréter une hausse du nombre de failles quand la couverture
          augmente.
\end{itemize}
\end{boxinfo}

% ============================================================================
\section*{Exercices}
% ============================================================================

\begin{boxexo}[title={\ding{45}\quad Exercice~11.1. Classer \niveaubase}]
Avec la politique de SOLITRANS, classez~: (a) faille inscrite au catalogue
KEV sur le portail des transitaires~; (b) faille de score 9,4 sur un serveur
interne, non exploitée~; (c) faille de score 7,2 sur un poste~; (d) faille
inscrite au catalogue KEV sur un serveur interne~; (e) faille de score 5,0
sur une imprimante.
\end{boxexo}

\begin{boxexo}[title={\ding{45}\quad Exercice~11.2. Le stock de P2 \niveaumoyen}]
\begin{enumerate}[(a)]
    \item Vérifiez le calcul de l'encadré « La politique de juin 2025 »~:
          combien de semaines pour absorber les P1 et les P2 à quarante
          corrections par semaine~?
    \item Combien de corrections par semaine faudrait-il pour traiter les
          211 P2 en 15 jours, après les 38 P1~?
    \item Proposez deux solutions autres que le renfort.
\end{enumerate}
\end{boxexo}

\begin{boxexo}[title={\ding{45}\quad Exercice~11.3. Instruire une exception \niveaumoyen}]
Le portail des transitaires présente une faille P1~; le correctif casse le
module de facturation, et l'éditeur annonce une version compatible dans six
semaines. Rédigez la demande d'exception complète, et dites qui doit la
signer.
\end{boxexo}

\begin{boxexo}[title={\ding{45}\quad Exercice~11.4. Lire les indicateurs \niveauexpert}]
Au T4 2025, le nombre total de failles ouvertes de SOLITRANS augmente de
12~\% alors que tous les indicateurs de délai s'améliorent.
\begin{enumerate}[(a)]
    \item Proposez deux explications compatibles avec l'encadré « SOLITRANS,
          trois trimestres ».
    \item Comment les départager~?
    \item Rédigez la phrase du rapport trimestriel qui présente ce chiffre au
          comité sans l'alarmer à tort ni le rassurer à tort.
\end{enumerate}
\end{boxexo}

\bigskip

\begin{boxinfo}
\textbf{Pour aller plus loin.}
\begin{itemize}
    \item Le NIST SP~800-40~\cite{nist80040}.
    \item L'ISO/IEC~30111~\cite{iso30111} et l'ISO/IEC~29147~\cite{iso29147},
          pour le traitement et la divulgation du côté des éditeurs.
\end{itemize}
\end{boxinfo}

\begin{boxinfo}
\textbf{Transition.} Le programme de SOLITRANS empêche les incidents~;
l'enquête de NYANGA en explique un. Le Chapitre~\ref{chap:artefact-vuln}
relie les deux~: comment remonter d'une trace à la faille qui l'a permise.
\end{boxinfo}
